\begin{lemma}
\label{lemma-radical-element}
Let $R$ be a Noetherian local normal domain of dimension $2$.
Let $\mathfrak p_1, \ldots, \mathfrak p_r$ be pairwise distinct
primes of height $1$. There exists a nonzero element
$f \in \mathfrak p_1 \cap \ldots \cap \mathfrak p_r$ such
that $R/fR$ is reduced.
\end{lemma}

\begin{proof}
Let $f \in \mathfrak p_1 \cap \ldots \cap \mathfrak p_r$ be a nonzero element.
We will modify $f$ slightly to obtain an element that generates a radical ideal.
The localization $R_\mathfrak p$ of $R$ at each height one prime
ideal $\mathfrak p$ is a discrete valuation ring, see
Algebra, Lemma \ref{algebra-lemma-characterize-dvr} or
Algebra, Lemma \ref{algebra-lemma-criterion-normal}.
We denote by $\text{ord}_\mathfrak p(f)$ the corresponding
valuation of $f$ in $R_{\mathfrak p}$. Let
$\mathfrak q_1, \ldots, \mathfrak q_s$
be the distinct height one prime ideals containing $f$.
Write $\text{ord}_{\mathfrak q_j}(f) = m_j \geq 1$ for each $j$.
Then we define $\text{div}(f) = \sum_{j = 1}^s m_j\mathfrak q_j$
as a formal linear combination of
height one primes with integer coefficients.
Note for later use that each of the primes $\mathfrak p_i$
occurs among the primes $\mathfrak q_j$.
The ring $R/fR$ is reduced if and only if
$m_j = 1$ for $j = 1, \ldots, s$. Namely, if $m_j$ is $1$ then
$(R/fR)\mathfrak q_j$ is reduced and
$R/fR \subset \prod (R/fR)_{\mathfrak q_j}$ as
$\mathfrak q_1, \ldots, \mathfrak q_j$ are the associated primes
of $R/fR$, see Algebra, Lemmas
\ref{algebra-lemma-zero-at-ass-zero} and
\ref{algebra-lemma-normal-domain-intersection-localizations-height-1}.

\medskip\noindent
Choose and fix $g$ and $N$ as in Lemma \ref{lemma-syspar}.
For a nonzero $y \in R$ denote $t(y)$ the number of primes minimal over $y$.
Since $R$ is a normal domain, these primes
are height one and correspond $1$-to-$1$ to the minimal primes of
$R/yR$ (Algebra, Lemmas \ref{algebra-lemma-minimal-over-1} and
\ref{algebra-lemma-normal-domain-intersection-localizations-height-1}).
For example $t(f) = s$ is the number
of primes $\mathfrak q_j$ occurring in $\text{div}(f)$.
Let $h \in \mathfrak m^N$. By Lemma \ref{lemma-minprimespoly} we have
\begin{align*}
t(f + h) & \leq \text{length}_{R/(f + h)}(R/(f + h, g)) \\
& = \text{length}_R(R/(f + h, g)) \\
& = \text{length}_R(R/(f, g))
\end{align*}
see Algebra, Lemma \ref{algebra-lemma-length-independent}
for the first equality.
Therefore we see that $t(f + h)$ is bounded independent of
$h \in \mathfrak m^N$.

\medskip\noindent
By the boundedness proved above we may pick
$h \in \mathfrak m^N \cap \mathfrak p_1 \cap \ldots \cap \mathfrak p_r$
such that $t(f + h)$ is maximal among such $h$. Set $f' = f + h$.
Given $h' \in \mathfrak m^N \cap \mathfrak p_1 \cap \ldots \cap \mathfrak p_r$
we see that the number $t(f' + h') \leq t(f + h)$.
Thus after replacing $f$ by $f'$ we may assume that for every
$h \in \mathfrak m^N \cap \mathfrak p_1 \cap \ldots \cap \mathfrak p_r$
we have $t(f + h) \leq s$.

\medskip\noindent
Next, assume that we can find an element $h \in \mathfrak m^N$ such that
for each $j$ we have $\text{ord}_{\mathfrak q_j}(h) \geq 1$ and
$\text{ord}_{\mathfrak q_j}(h) = 1 \Leftrightarrow m_j > 1$.
Observe that
$h \in \mathfrak m^N \cap \mathfrak p_1 \cap \ldots \cap \mathfrak p_r$.
Then $\text{ord}_{\mathfrak q_j}(f + h) = 1$
for every $j$ by elementary properties of valuations.
Thus
$$
\text{div}(f + h) = \sum\nolimits_{j = 1}^s \mathfrak q_j +
\sum\nolimits_{k = 1}^v e_k \mathfrak r_k
$$
for some pairwise distinct height one prime ideals
$\mathfrak r_1, \ldots, \mathfrak r_v$ and $e_k \geq 1$.
However, since $s = t(f) \geq t(f + h)$ we see that $v = 0$
and we have found the desired element.

\medskip\noindent
Now we will pick $h$ that satisfies the above criteria.
By prime avoidance (Algebra, Lemma \ref{algebra-lemma-silly})
for each $1 \leq j \leq s$ we can find an element $a_j \in \mathfrak q_j$
such that $a_j \not \in \mathfrak q_{j'}$ for $j' \not = j$
and $a_j \not \in \mathfrak q_j^{(2)}$. Here
$\mathfrak q_j^{(2)} = \{x \in R \mid \text{ord}_{\mathfrak q_j}(x) \geq 2\}$
is the second symbolic power of $\mathfrak q_j$.
Then we take
$$
h = \prod\nolimits_{m_j = 1} a_j^2 \times
\prod\nolimits_{m_j > 1} a_j
$$
Then $h$ clearly satisfies the conditions on valuations imposed above.
If $h \not \in \mathfrak m^N$, then we multiply by an element of
$\mathfrak m^N$ which is not contained in $\mathfrak q_j$ for all $j$.
\end{proof}